\begin{table}[H]
\centering
\captionsetup{justification=centering}
\caption{Resultados pontuais do modelo multivariado para a taxa de desocupação - 4\textordmasculine{} trimestre de 2024}
\label{tab:pontual_taxa}

\renewcommand{\arraystretch}{0.8}
\resizebox{\linewidth}{!}{%
\begin{tabular}{@{}p{2.8cm}cccccc@{}}
\toprule
\textbf{Estrato Geográfico} & $\hat{y}$ & $CV(\hat{y})$ & IC 95\% & $\hat{\theta}$ & $CV(\hat{\theta})$ & IC 95\% \\
\midrule
01 - Belo Horizonte & 4,59 & 13,92\% & [3,34 ; 5,84] & 5,37 & 6,82\% & [4,66 ; 6,09] \\
\raisebox{0em}{\parbox[t]{2.8cm}{02 - Entorno e Colar \\[-2pt] Metropolitano de BH}} & 5,16 & 10,16\% & [4,13 ; 6,18] & 4,69 & 8,29\% & [3,93 ; 5,45] \\
03 - Sul de Minas & 2,61 & 16,25\% & [1,78 ; 3,44] & 2,82 & 10,53\% & [2,24 ; 3,40] \\
04 - Triângulo Mineiro & 3,69 & 16,16\% & [2,52 ; 4,85] & 3,57 & 8,24\% & [3,00 ; 4,15] \\
05 - Zona da Mata & 4,31 & 13,05\% & [3,21 ; 5,41] & 4,26 & 10,39\% & [3,40 ; 5,13] \\
06 - Norte de Minas & 5,63 & 16,14\% & [3,85 ; 7,41] & 5,77 & 9,08\% & [4,74 ; 6,79] \\
07 - Vale do Rio Doce & 3,98 & 13,39\% & [2,94 ; 5,03] & 3,88 & 10,63\% & [3,07 ; 4,69] \\
08 - Central & 4,07 & 15,09\% & [2,86 ; 5,27] & 3,53 & 10,34\% & [2,82 ; 4,25] \\
\bottomrule
\end{tabular}%
}
\fonte{Elaboração própria, com base nos dados da PNAD Contínua. Nota: \(\hat{y}\) é a estimativa direta e \(\hat{\theta}\) o sinal estimado pelo modelo multivariado (tendência mais sazonalidade). Valores em pontos percentuais.}
\end{table}
